\section{Related Work and Novelty Boundary}
\label{sec:related}
\textbf{Structure and detection.} Zimmermann et al. study npm dependency and maintainer concentration; Zahan et al. examine metadata weak links~\cite{zimmermann2019,zahan2022}. These establish the feature families rather than their novelty here. Ohm et al. organize documented supply-chain attacks~\cite{ohm2020}. Amalfi and a recent detector benchmark evaluate released package content~\cite{sejfia2022,guo2026}. A benign release-classification label does not establish that its publisher remains uncompromised in a future horizon.

\textbf{Workflows and authority.} CI workflow characterization and ARGUS analyze security properties and injection paths beyond directory presence~\cite{koishybayev2022,muralee2023}. GitHub documents permission and privileged-trigger semantics~\cite{githubsecure}. Release-authority measurement compares predecessor release paths and reports no exact malicious-version overlap with its policy triggers in the studied cohort~\cite{santos2026}. Authority discontinuity is consequently a distinct review surface.

\textbf{Historical data and forecasting.} npm-follower preserves releases that can later disappear, and CHRONO-RESOLUTION provides release-point dependency resolution~\cite{pinckney2023,rahman2026}. SupplyTGN evaluates temporal vulnerability-risk ranking in a PyPI/OSV setting; PatchRisk studies future transitive vulnerability exposure for npm and PyPI~\cite{supplytgn2026,patchrisk2026}. These narrow any broad claim to introduce temporal package forecasting.

Table~\ref{tab:positioning} locates our empirical contribution: disagreement between co-located artifacts and its queue consequences. We claim neither a new graph architecture nor the first temporal package benchmark. The literature screen, updated on 20 October 2026, is targeted rather than systematic.
\begin{table*}[t]
\caption{Positioning against the closest research streams. Differences identify scope, not superiority.}
\label{tab:positioning}
\centering\footnotesize
\begin{tabularx}{\textwidth}{@{}>{\raggedright\arraybackslash}p{0.22\textwidth}YYY@{}}
\toprule
Work & Principal evidence & Evaluated target or artifact & Boundary for this study \\
\midrule
Small World; Weak Links~\cite{zimmermann2019,zahan2022} & Dependencies, maintainers, metadata & Concentration and metadata weaknesses & Structural descriptors already established \\
Amalfi; detector benchmark~\cite{sejfia2022,guo2026} & Released package content & Malicious-package detection & Detection labels do not define future non-events \\
CI characterization; ARGUS~\cite{koishybayev2022,muralee2023} & Workflow and action code & Configuration properties; injection paths & Directory presence is a weaker observation \\
On Good Authority~\cite{santos2026} & Predecessor release-path evidence & Authority discontinuities & Authority change and compromise differ \\
npm-follower; CHRONO-RESOLUTION~\cite{pinckney2023,rahman2026} & Archived releases; release-point resolution & Historical data resources & Graph reconstruction alone is not new \\
SupplyTGN; PatchRisk~\cite{supplytgn2026,patchrisk2026} & Temporal package/advisory evidence & Future vulnerability-risk ranking & Temporal forecasting is already studied \\
This study & Two local archives; incident-record inventory & Metric consistency, queue drift, feasibility & No fitted model or predictive benchmark result \\
\bottomrule
\end{tabularx}
\end{table*}

